% !TEX root = ../main.tex
% OWNER: Matthew | ~300 words

We split the period from 1 November 2025 to 9 October 2026 into
consecutive 14-day windows (fortnights) of 336 hours each. Test windows
start on 28 February 2026, the first day of the shock, and eight baseline
windows run back from it. The model is frozen throughout, so any change
in its explanations comes from the data and not from retraining.

\textbf{Explanation drift.} For each window we compute TreeSHAP values
\cite{lundberg2020} for every hour and score each feature by its mean
absolute SHAP value over the window, which ranks the features by how
strongly they moved the forecast. The reference ranking comes from the
October 2025 validation month. We compare each window's ranking to the
reference with weighted Kendall's tau \cite{kendall1938,vigna2015}
\begin{equation}
\tau_w = \frac{\sum_{i<j} w_{ij}\,\mathrm{sgn}(x_i - x_j)\,\mathrm{sgn}(y_i - y_j)}
              {\sqrt{\sum_{i<j} w_{ij}\,\mathrm{sgn}(x_i - x_j)^2}\,
               \sqrt{\sum_{i<j} w_{ij}\,\mathrm{sgn}(y_i - y_j)^2}},
\end{equation}
where $x_i$ and $y_i$ are feature $i$'s scores in the two windows. Each
pair is weighted by $w_{ij} = \frac{1}{r_i+1} + \frac{1}{r_j+1}$, where
$r$ is a feature's rank (0 for the top), so pairs near the top of the
ranking count most \cite{vigna2015}. We weight because a few features
carry most of the SHAP mass. In October the top four of twelve account
for 79\% of it, so a swap between the top two matters far more than one
between the bottom two. The resulting $\tau_w$ measures how far a
window's ranking has moved from the October reference. A value of 1
means the same order, 0 means no relation and $-1$ means a fully
reversed order. As a seasonal check, we also compare each window with
the same 14 days one year earlier, shifted 364 days so weekdays align.

\textbf{Input drift.} As a check that ignores the model, we compute the
Population Stability Index (PSI) \cite{yurdakul2018} for each input
except hour, weekday, month and the holiday flag, whose values are set
by the calendar and cannot drift. PSI measures how much the spread of
an input's values has shifted compared with the same 14 days one year
earlier, where 0 means no shift. We average the eight inputs' PSI values
to give one PSI per window, and also flag any single input at or above
0.25, a commonly used cut off \cite{yurdakul2018}.

\textbf{Warning rule.} We treat the eight baseline windows as normal
behaviour. For each of $\tau_w$, MAE and PSI we take its mean over these
windows, which gives its usual level, and its standard deviation (SD),
which gives how much it usually varies. A test window raises a warning
when it lies more than 2~SD from the mean in the worse direction, which
is lower for $\tau_w$ and higher for MAE and PSI. For the blind model
this puts the $\tau_w$ warning line at 0.71. Because MAE gets its own
2~SD warning line in the same way, we can measure lead time as the
number of days between the first $\tau_w$ warning and the first MAE
warning. We also report a stricter variant that needs two consecutive
windows over the line.